\chapter{Exact Agent-Network Coarse-Graining}
\label{comp:ch11}

\section{The question and the answer}
\label{comp:ch11-question}

The technical chapter (\MT{ch:agent-network-rg}) applies one coarse channel, a Markov kernel from
fine to coarse configurations (\MT{def:cg-coarse-channel}), to the interacting agent network at the
base point $c_*$ and asks what the step does to the network's energy function and couplings, and
what must be carried along so that nothing is lost. The coarse effective action is minus the
logarithm of a conditional partition function: the fine Boltzmann weight averaged over the fine
configurations compatible with each coarse configuration (\MT{thm:rg-effective-action}). Written in
interaction coordinates, one component for every set of coupled agents, the step is an exact
nonlinear map from fine couplings to coarse couplings, and the derivative of any composite of such
steps along an exact trajectory has norm at most $3^{\lvert V_n\rvert}-1$, a constant fixed by the number of target agents and
independent of the number of steps (\MT{prop:rg-interaction-native-growth-bound}). What must be
kept is collected in a table (\MT{tab:rg-retained-state}): induced multi-agent couplings, the joint
law of fine and coarse variables, the joint law of attention events rather than attention rows, the
gauge marks of each block, and the path law with its memory.

Has this resolved the renormalization of the agent network at one base point? Not in the sense a
physicist gives the word, and the chapter does not say that it has. Exact and proved are the
effective action as minus the log density of the pushed evidence measure relative to the pushed
reference, for every channel; the
exact step on couplings, together with an identity for the residual of any declared truncation; the
finite-step multipliers (the linearized exact flow) and their depth-independent bound; and the
coarse path law with an exact memory recurrence, filter, and finite-memory cost. Absent are four
things a physicist would expect. There is no closed truncation: apart from sectors produced by
independence, symmetry, or couplings already carried by the coarse variables, no proper subspace of
couplings is shown to be mapped into itself, and no truncation residual is estimated
(\MT{open:rg-invariant-retained-subspace}). There is no flow on one fixed set of couplings: across a
step that reduces the number of configurations the coupling spaces have different dimensions and
admit no identification (\MT{prop:rg-interaction-finite-network-obstruction}), and the main
manuscript constructs no rescaling, fixed point, or beta function (\MT{sec:intro-conventions}).
There is no thermodynamic limit: every statement concerns a finite network, and infinite networks
are deferred to Supplement B, where the limit questions are open (\SB{open:rg-two-index-limits}).
And there is no parent dynamics: the parent exists at every instant as an exact pushforward
(\MT{thm:rg-pointwise-parent-datum}), but whether its evolution closes on itself is
\MT{open:rg-pointwise-parent-dynamics}, the defining open problem of the current phase.

The chapter uses the exact coarse free energy of \cref{comp:ch09} (\MT{thm:rg-exact-coarse-vfe}) and
the root and holonomy marks of \cref{comp:ch10}, and it feeds \cref{comp:ch12} and the formation
witnesses of \cref{comp:ch13}, whose parents are exact marked parents in the present sense.

\section{The picture}
\label{comp:ch11-picture}

Take $N$ agents at $c_*$ with states $y_i$ (belief samples, model presentations, labels) and a joint
law equal to a reference probability times a record likelihood. Read the reference $\varrho$ as the a
priori measure, $L_o=e^{-H_o}$ as a Boltzmann weight at unit temperature, $\mathsf m_o=e^{-H_o}\varrho$
as the Gibbs weight, and its mass $M_o$ as the partition function; in Bayesian terms $M_o$ is the
evidence and the normalized $\bar{\mathsf m}_o$ the posterior. A Kadanoff block map is a
deterministic channel sending each block to one block variable; a noisy block map is a stochastic
channel; both are admitted, and random or soft blocks enter through the structural channel of
\cref{comp:ch12}. The exact coarse theory is the pushforward of every law. The chapter says what
that pushforward looks like as an effective Hamiltonian with couplings, and what a reduced
description cannot discard.

\begin{analogy}{exact}
The conditional-partition formula is the position-space renormalization formula with a weight
function, in the block-spin tradition of \citet{kadanoff1966scaling} reviewed by
\citet{wilson1974renormalization}: $e^{-H'(s')}=\sum_sP(s'\given s)e^{-H(s)}$ up to normalization.
The a priori measure is explicit, so fiber-size (entropy) terms such as those of a majority rule are
carried by the pushed reference. Lengths and fields are not rescaled and nothing is truncated; the
identification is exact at the level of the formula and brings no fixed-point theory with it.
\end{analogy}

\begin{analogy}{exact}
The two-agent example of \MT{rem:rg-two-agent-example} is decimation: summing over spin $2$ of a
pair with fields $\xi_1,\xi_2$ and coupling $\xi_{12}$ gives the coarse field
$\xi_1'=\xi_1+\operatorname{artanh}(\tanh\xi_2\tanh\xi_{12})$, so the pair coupling survives only as a
field and only through the polarization $\tanh\xi_2$ of the eliminated spin. With three neighbors the
same mechanism produces a three-spin coupling (\MT{prop:rg-hyperedge-generation}), as decimation of
the two-dimensional Ising model produces four-spin couplings. This is why exact coarse-graining
leaves the pairwise world.
\end{analogy}

\section{Objects and notation}
\label{comp:ch11-notation}

\begin{decoder}
depth $\ell$, $V_\ell$ & Coarse-graining level and its finite agent set & Blocking step; not time\\
$C_\ell$, $C$ & Coarse channel & Block-spin weight function $P(s'\given s)$\\
$\varrho$, $\mathsf m_o$, $M_o$ & Reference probability; evidence measure $e^{-H_o}\varrho$; its mass
& A priori measure; Gibbs weight; partition function\\
$H_o$, $H^c_o$ & Action (minus log density relative to the reference); coarse action & Hamiltonian
at unit temperature; block Hamiltonian\\
$C^\dagger$, $\mathsf U$ & Reverse (Bayes) kernel of fine given coarse; its conditional expectation &
Conditional Gibbs measure at fixed block spins\\
$\mathcal A(\psi)$, $\mathsf U^\psi$ & Exact change of the coarse action under a bounded change $\psi$
of the fine action; tilted reverse expectation & Change of the effective Hamiltonian; perturbed
conditional ensemble\\
$\lVert\cdot\rVert_{\mathrm{osc}}$, $\delta$ & Half the range of a function; largest total-variation
distance between reverse conditionals & Sup norm of an energy modulo constants; Dobrushin
coefficient\\
$\nu_\ell$, $\mathsf P_{\ell,B}$ & Product reference law; projector onto the pure $B$-body part &
Noninteracting a priori measure; projection onto $k$-body terms\\
$\mathcal G_\ell$, $\mathsf g=(\mathsf g_B)$ & One bounded component per nonempty set of agents,
normed by $\sum_B\lVert\mathsf g_B\rVert_\infty$ & All $k$-body couplings of a finite system\\
$\operatorname{asm}$, $\operatorname{hd}$ & Couplings to action; action to couplings & Writing and
reading off a Hamiltonian\\
$\mathcal T_\ell$ & Exact interaction step & Exact coupling recursion, untruncated, unrescaled\\
$D_\ell$, $D_{n\leftarrow\ell}$ & Derivative of one step on an exact orbit; ordered product &
Linearized RG matrix along a trajectory\\
$\mathsf P^{\mathrm{ret}}_\ell$, $r^{\mathcal G}$ & Declared retained projection; its residual &
Truncation of coupling space; truncation error\\
hyperedge & Component on three or more agents & Multi-spin coupling\\
posterior bridge, bridge kernels & Joint posterior law of fine and coarse variables; its pair
conditionals & Joint law of spins and block spins\\
$\eta$, $\alpha$, $\beta$ & Event law over edge copies; receiver occupancy; attention row & Who
attends to whom; its receiver marginal; the conditional\\
$D^{(k)}_N$, $\mathrm h^{\mathrm{ex}}_k$ & Cost of a memory-$k$ path law; excess conditional entropy
& Truncation cost of a $k$-th order Markov model; excess entropy\\
\end{decoder}

Some glyphs mislead. $M_o$ is an evidence, not a magnetization, and $\mathcal G_\ell$ is a coupling
space, not a group. Three objects share a $D$: the multipliers $D_\ell$, the edge divergences $D_e$,
and the memory defects $D^{(k)}_N$. $H$ is an action early in the chapter and an entropy in the memory
theorems. The dagger in $C^\dagger$ is Bayes reversal (the $L^2$ adjoint of the lift), not complex
conjugation. Here $z$ is the coarse sample of one arrow, whereas in \cref{comp:ch12} $Z_\ell$ is the
coordinate block at depth $\ell$. ``Exact'' means without approximation, not solvable in closed form;
``closure'' means that the kept state suffices for the next exact step; a ``multiplier'' is a
derivative and a ``cocycle'' an ordered product of derivatives along an orbit, not a cohomology
class. The main manuscript avoids the words relevant, fixed point, and beta function by design.

\section{The formulas, explained}
\label{comp:ch11-formulas}

\subsection{The effective action is a conditional partition function}

For one arrow $C:\mathsf Y\rightsquigarrow\mathsf Z$, with $C^\dagger_\varrho(z,dy)$ the law of the
fine state given the coarse state under the reference joint law $\varrho(dy)C(y,dz)$,
\[
\exp[-H^c_o(z)]=\int\exp[-H_o(y)]C^\dagger_\varrho(z,dy),
\qquad
M_o=\int\exp[-H^c_o(z)]\varrho^c(dz),
\qquad\varrho^c=\varrho C
\]
(\MT{thm:rg-effective-action}). The coarse Boltzmann weight at $z$ is the a priori average of the
fine weight over the configurations the channel can send to $z$, and the partition function is the
same at both scales. For the identity channel $H^c_o=H_o$; for a channel to a point the action
becomes the number $-\log M_o$; for a deterministic block map on a finite space with uniform
reference, $e^{-H^c_o(z)}$ is the average of $e^{-H_o}$ over the fiber of $z$, the Kadanoff block
Hamiltonian with the entropy $\log\lvert\chi^{-1}(z)\rvert$ moved into the reference. The theorem holds for any channel between standard Borel spaces
and is standard in substance (the density of a pushforward is a conditional expectation). It says
nothing about the form of $H^c_o$: on a finite network it always exists but may couple every coarse
agent to every other, while in infinite volume the analogous object can fail to exist as a
quasilocal Hamiltonian \citep{VanEnterFernandezSokal1993}, which no finite-network statement reaches.

\subsection{Perturbing the fine action: cumulants and contraction}

Perturb the fine evidence measure by a bounded $\psi$. With $C^\dagger$ the reverse kernel of the
normalized joint law,
\[
\mathcal A(\psi)(z)=-\log\int e^{-\psi(y)}C^\dagger(z,dy),
\qquad
D\mathcal A(\psi)[h]=\mathsf U^\psi h,
\qquad
D^2\mathcal A(\psi)[h,k]=-\operatorname{Cov}_{C^\dagger_\psi(z,\cdot)}(h,k)
\]
(\MT{thm:rg-bounded-action-calculus}). $\mathcal A$ is the conditional cumulant generating function of
$\psi$, $\langle\psi\rangle_z-\tfrac12\langle\psi;\psi\rangle_z+\cdots$, the cumulant expansion of any
integrating-out calculation. It is real analytic on all bounded perturbations, order preserving,
passes constants through, and never expands the sup norm
(\MTeq{eq:rg-nonlinear-action-sup-contraction}). Its linear part $\mathsf U$ is a conditional
expectation: it contracts every $L^p$ norm, its $L^2$ loss is the expected conditional variance, and
it contracts oscillation by the Dobrushin coefficient $\delta$ of the reverse kernel
\citep{Dobrushin1956} (\MT{thm:rg-action-lp-contraction}). For Kadanoff blocks the last bound is
empty: under a deterministic block map with a nondegenerate coarse law the reverse conditionals at
distinct coarse states live on disjoint fibers, so $\delta=1$ and the bound says only that nothing
expands; strict contraction needs a noisy channel. Every norm here is a norm of the whole action. A
uniform perturbation of size $h$ per agent has sup norm of order $Nh$, so these statements control
the total action, not couplings per agent.

\subsection{Interaction coordinates}

Under a product reference $\nu_\ell=\bigotimes_i\nu_{\ell i}$ equivalent to the normalized law
(\MT{def:rg-network-scale}), let $\E_{\ell,B}$ average out the agents outside $B$. The Hoeffding
projectors and the assembly of an action from couplings are
\[
\mathsf P_{\ell,B}=\sum_{B'\subseteq B}(-1)^{\lvert B\rvert-\lvert B'\rvert}\E_{\ell,B'},
\qquad
\psi_{\mathsf g}=\sum_{\varnothing\ne B\subseteq V_\ell}\mathsf g_B
\]
(\MT{def:rg-hoeffding-coordinates}). $\mathsf P_{\ell,B}f$ is the part of $f$ that depends on all of
$y_B$ jointly and has zero average in each single agent of $B$: the cluster, or ANOVA, expansion of the
Hamiltonian into $k$-body terms. For $\pm1$ spins with uniform reference the $B$-component is
$J_B\prod_{i\in B}y_i$. Gibbs-measure theory also uses the vacuum normalization relative to a reference
configuration \citep{Georgii1988}; the chapter has it too (\MTeq{eq:rg-mobius-potentials}) but uses only
the Hoeffding components in the step. \MT{thm:rg-hoeffding-action-isomorphism} proves that couplings
and actions modulo constants are isomorphic, with $\lVert\operatorname{asm}_\ell\rVert\le1$ and
$\lVert\operatorname{hd}_\ell\rVert\le3^{\lvert V_\ell\rvert}-1$, sharp: for independent biased $\pm1$
coordinates the product $\prod_iY_i$, of oscillation one, has components whose norms sum to nearly
$3^n-1$. The size of an action and the total size of its couplings are equivalent only up to a factor
exponential in the number of agents. The decomposition is gauge covariant when the frame change acts
agent by agent (\MT{cor:rg-hoeffding-gauge-covariance}), which is the case for the program's vertex
gauge; preserving the product reference alone is not enough.

Two premises limit the scope, and both matter for the program. The bounded tier
(\MT{rem:rg-bounded-tier-scope}) admits finite and categorical agents and continuous couplings of
compact group-valued states, so $\mathrm{SU}(N)$-valued states are inside; it excludes every polynomial
coupling on $\R^K$, since $y_i^\top\mathsf Ny_j$ is unbounded. The Gaussian realization with $\GL(K,\R)$
or $\GL(K,\mathbb C)$ links is therefore outside this machinery; its exact marginalization (Kron
reduction, a Schur complement) is in \SG{ch:gaussian-coarsegraining}. And product equivalence at both
depths, automatic for strictly positive laws on finite products, is not preserved by every channel: a
channel cloning one binary agent into two yields a diagonal law that no product law matches
(\MT{prop:rg-product-equivalence-not-preserved}).

\subsection{The exact step and the native growth bound}

The exact interaction step is
$\mathcal T_\ell=\operatorname{hd}_{\ell+1}\overline{\mathcal A}_\ell\operatorname{asm}_\ell$
(\MT{def:rg-interaction-step}): assemble the couplings into an action, integrate out exactly modulo
constants, read off the coarse couplings. The couplings perturb a base law whose own image is carried
separately, so $\mathcal T_\ell(0)=0$ for every model and channel (\MT{lem:rg-interaction-step}). Along an exact orbit the multipliers are
$D_\ell=\operatorname{hd}_{\ell+1}\overline{\mathsf U^{\psi_{\mathsf g_\ell}}_\ell}\operatorname{asm}_\ell$,
every interior pair $\operatorname{asm}_k\operatorname{hd}_k$ cancels, and
\[
D_{n\leftarrow\ell}
=\operatorname{hd}_n\overline{\mathsf U_{n-1}^{\psi_{\mathsf g_{n-1}}}}\cdots
\overline{\mathsf U_\ell^{\psi_{\mathsf g_\ell}}}\operatorname{asm}_\ell,
\qquad
\lVert D_{n\leftarrow\ell}\rVert_{\mathcal G_\ell\to\mathcal G_n}\le3^{\lvert V_n\rvert}-1
\]
(\MT{prop:rg-interaction-native-growth-bound}). Every tilted conditional expectation has oscillation
norm at most one, assembly has norm at most one, and only the final extraction can exceed one.

What is bounded is the operator norm of the linearized composite step from all depth-$\ell$ couplings
to all depth-$n$ couplings, each measured by the sum of the sup norms of its components: a
perturbation of total size one in the fine couplings changes the total size of the depth-$n$ couplings
by at most $3^{\lvert V_n\rvert}-1$ to first order, however many steps intervene and however large the
fine network was. It is not a bound per coupling or per agent. As a number it is useful only near the
top of a blocking: with pair blocks $512$ agents admit nine steps, and the constant is $2$, $8$, $80$,
and $6560$ for targets of one, two, four, and eight blocks, but about $10^{122}$ for the first step's
target of $256$ agents. The real content is independent of $N$: in the action norm the linearized
exact flow is a product of contractions, and all apparent amplification is one change of norm at the
target. Supplement B turns this into the statement that the exact step has no relevant direction under
tempered comparisons (\SB{cor:rg-interaction-no-relevant-mode}); relevance requires an extensive
normalization declared from outside.

\begin{analogy}{structural}
In Wilson's renormalization a uniform field $h\sum_is_i$ is relevant because the per-site field grows
as $b^{y_h}h$ under blocking by $b$ in $d$ dimensions, while the whole-action size $hN$ becomes
$b^{y_h-d}hN$, which does not grow since $y_h\le d$. Relevance is a statement about couplings per
site. The native bound is the exact finite-network form of the other half: with no rescaling and no
per-site normalization, nothing grows.
\end{analogy}

\subsection{Residuals, invariant sectors, and the finite-network obstruction}

A declared truncation is a gauge-intertwining idempotent $\mathsf P^{\mathrm{ret}}$ on couplings, for
example onto single-agent and pair components (\MT{def:rg-retained-projection}), with residual
$r^{\mathcal G}_{\ell+1}=(I-\mathsf P^{\mathrm{ret}}_{\ell+1})\mathcal T_\ell(\mathsf g_\ell)$.
\MT{prop:rg-interaction-residual} proves that the coupling and action residuals vanish together, agree
in norm up to $3^{\lvert V_{\ell+1}\rvert}-1$, and vanish on the retained range exactly when that range
is mapped into its successor. It is an identity, not an estimate. The proved invariant pairs are
generated by structure, not dynamics (\MT{prop:rg-interaction-intra-block-invariance},
\MT{rem:rg-trivial-invariant-pairs}): independent blocks under a block-product channel generate no
inter-block couplings; a finite symmetry preserves its fixed sector, such as the even sector under a
global spin flip; and couplings already carried by the coarse variables of a deterministic channel
pass through unchanged. Anything beyond is \MT{open:rg-invariant-retained-subspace}.

For finite agent spaces the coupling space has dimension $\prod_i\lvert\operatorname{supp}\nu_{\ell
i}\rvert-1$, so $2^N-1$ for $N$ binary agents (\MT{prop:rg-interaction-finite-network-obstruction}). A
step that reduces the number of configurations reduces the dimension (six binary agents in three binary
blocks go from $63$ to $7$), so no identification of the two coupling spaces exists and no single space
carries a flow; if every step divides the vertex count by at least $b>1$, the depth is at most
$\log_b\lvert V_\ell\rvert$. This is not a small-network artifact. It holds at $N=512$ and at
$N=6\times10^{23}$ alike: a finite system can be blocked only logarithmically often, and its coupling
space shrinks with each blocking. Physics removes the obstruction by passing to infinite volume and
identifying couplings of the same shape across scales by translation invariance; neither ingredient is
in the main manuscript, and Supplement B records that finite networks do not carry the asymptotic
classification (\SB{prop:rg-finite-depth-no-asymptotics}). The dimension count needs finite agent
spaces; termination holds for all.

\subsection{Hyperedges and bridge kernels}

Marginalizing a variable multiplies the factors that touch it and leaves one factor on the union of its
neighbors, so iterated elimination produces factors on elimination cliques
(\MT{prop:rg-hyperedge-generation}); this is variable elimination with fill-in. Eliminating the center
$y_0$ of a star of three $\pm1$ neighbors, with field $\xi_0$ and couplings $\xi_r$, leaves a
coefficient $2\operatorname{sech}^2(\xi_0)\tanh(\xi_0)\xi_1\xi_2\xi_3+O(\xi^5)$ on $y_1y_2y_3$: nonzero
for small nonzero couplings when the eliminated spin carries a field, and identically zero at
$\xi_0=0$, where spin-flip symmetry kills odd terms. Pairwise closure fails in general; ``generic'' is
shown within the star family, not over all networks, although the mechanism is that of decimation.

The posterior bridge is the joint posterior law of fine and coarse variables, and the bridge kernels
are its pair conditionals between a fine agent and a meta-agent, Bayes adjoints of each other
(\MT{def:rg-posterior-bridge}, \MT{prop:rg-bridge-kernels}). The meta-to-agent kernel does not
reconstruct the fine posterior: when the likelihood couples blocks, a fine agent's conditional law
depends on neighboring block variables too, just as the configuration inside a block is not fixed by
its own block spin. For a deterministic block map, lifting and conditioning returns a coarse observable
unchanged; for a randomized channel it does not.

\subsection{Meta-attention}

An attention row is the conditional law of a source given a receiver, and a conditional law cannot be
coarse-grained without the law of what it conditions on, as lumping a transition matrix needs
occupation weights. The recursive state is the joint event law $\eta$ over oriented edge copies:
\[
\eta^c_{IJ}(z)=\E_{\mathbb B_o}\Bigl[\sum_{e:j\to i,\ i\in I,\ j\in J}\eta_e(Y)\Bigm|Z=z\Bigr],
\qquad
\alpha^c_I=\sum_J\eta^c_{IJ},
\qquad
\beta^c_{IJ}=\eta^c_{IJ}/\alpha^c_I
\]
(\MT{prop:rg-meta-attention-pushforward}). The next depth pushes $\eta^c$, not $\beta^c$, and events
inside a block become meta-self events, so a meta-agent attends to itself whenever its members attend
to each other.

Where the receiver weights come from is a modeling choice with consequences. In the single-edge model
(\MT{hyp:rg-single-edge-attention}) each record carries one joint receiver--source label, and the
receiver law within a block is evidence weighted,
$\alpha_i'\propto a_{i\mid I}\widetilde{\mathrm c}_i\mathcal Z_i$, where
$\mathcal Z_i=\sum_j\pi_{ij}e^{-D_{ij}/\tau_i}$ is the partition function of receiver $i$'s row and
$\widetilde{\mathrm c}_i$ the likelihood ratio of its record under the event; the block row is the
$\alpha'$-weighted average of member rows (\MT{prop:rg-attention-single-edge-weights}). In the
per-receiver model of \cref{comp:ch06}, every receiver has its own label, and a block row needs a
structural selector that draws one receiver with fixed weights independently of everything; its
posterior equals its prior, and the block row collapses to the prior-weighted average
(\MT{prop:rg-attention-per-receiver-collapse}). The two agree when $\widetilde{\mathrm c}_i\mathcal Z_i$
is constant on the block and otherwise differ ($3/4$ against $1/2$ in the chapter's two-receiver
example); the collapse is the loss of evidence weighting. Keeping all receivers' labels is exact, and
collapsing them to one row is a further channel whose cost enters \MTeq{eq:rg-vfe-chain-rule}.

\begin{analogy}{exact}
Grouping the sources of one row gives the Gibbs row of the coarse prior $\pi^c_J=\sum_{j\in J}\pi_j$ and
coarse energy $D^c_J=-\tau\log[(\pi^c_J)^{-1}\sum_{j\in J}\pi_je^{-D_j/\tau}]$
(\MT{rem:rg-attention-log-sum-exp}): the free energy of a macrostate whose microstates have
degeneracies $\pi_j$. In the single-edge model, receivers are merged with weights proportional to their
own partition functions, as subsystems are when their ensembles are merged.
\end{analogy}

The event-weighted mean of transported link representations between blocks maps the root frame of $J$
to that of $I$ covariantly, but it need not lie in the group (the mean of rotations by $\pm\pi/2$ is
zero) or be invertible, and it is only a first moment of the joint law of marks and features that
exact closure keeps (\MT{prop:rg-attention-hom-moment}). It is the blocked link of lattice gauge
theory, which leaves the group and is usually projected back; here nothing is projected.

\subsection{Path laws and memory}

Applying the channel at every time of a path is again a coarse channel, so the free-energy chain rule
holds for path laws; the ordering index is an inference index, not physical time
(\MT{sec:rg-path-space}). The coarse path law is exact but generally not Markov.
\MT{thm:rg-strong-lumpability} is the exact condition of \citet[Thm.~6.3.2]{KemenySnell1976}, on
standard Borel spaces, for Markov coarse dynamics from every initial law: equal jump probabilities into
each coarse set from all fine states with the same coarse value. \MT{rem:rg-lumpability-witnesses}
shows a chain whose coarse process is Markov from one start without lumpability, and a deterministic
three-cycle with labels $(0,0,1)$ whose stationary coarse process has memory. For a linear recursion
$\xi_{n+1}=\mathsf T\xi_n$ with resolved coordinate $w_n=\mathsf C\xi_n$, lift $\mathsf J$
($\mathsf C\mathsf J=I$), and $\mathsf Q=I-\mathsf J\mathsf C$,
\[
w_{n+1}=\mathsf C\mathsf T\mathsf Jw_n
+\sum_{k=0}^{n-1}\mathsf C\mathsf T\mathsf Q(\mathsf Q\mathsf T\mathsf Q)^{n-1-k}\mathsf Q\mathsf T\mathsf Jw_k
+\mathsf C\mathsf T\mathsf Q(\mathsf Q\mathsf T\mathsf Q)^n\mathsf Q\xi_0
\]
exactly (\MT{thm:rg-projection-memory}): a Markovian term, memory kernels, and the unresolved initial
state. This is the discrete, linear form of the projection-operator method of \citet{Nakajima1958} and
\citet{Zwanzig1960}, the last term playing the role of the noise.

For a hidden finite chain observed through labels, the exact filter is the forward algorithm, and its
terminal posterior determines every prediction but may take infinitely many values
(\MT{thm:rg-memory-filter}), so a finite-state Markov parent does not follow even in principle. The
memory-$k$ approximation whose rows are the true conditionals given the last $k$ symbols is the unique
forward-KL optimum, with cost
\[
D^{(k)}_N=\KL\bigl(\mathbb P^{\mathrm{path}}_N\Vert\widehat{\mathbb P}^{(k)}_N\bigr)
=\sum_{t=0}^{N-1}I\bigl(Z_t;Z_{0:t-k-1}\given Z_{\max(0,t-k):t-1}\bigr)
\]
(\MT{thm:rg-memory-projection}): what the discarded past still says about the present given the kept
past. The reverse divergence is not controlled (the memory-one projection of the three-cycle gives the
impossible word $000$ probability $1/6$), and using the projection as a prior changes a recognition
law's free energy by an unsigned amount. For a stationary start,
\MT{thm:rg-memory-duration} gives $D^{(k)}_N/N\to\mathrm h^{\mathrm{ex}}_k$ with
$\sum_j\mathrm h^{\mathrm{ex}}_j\le H(Y_0)$ and $\mathrm h^{\mathrm{ex}}_k\le H(Y_0)/(k+1)$, where $N$ is
now the path length: the excess entropy \citep{CrutchfieldFeldman2003} is bounded by the entropy of the
hidden state. For the program this budget is extensive, up to $512\log2\approx355$ nats for $512$
binary agents, and the theorem is stated for finite hidden spaces only. The prepared parent of
\cref{comp:ch08} is Markov when the fine state inside each parent is refreshed to its exact conditional
law at every step (\MT{thm:rg-memory-prepared-parent}), which is adiabatic elimination by hypothesis
(\MT{hyp:rg-memory-instantaneous-refresh}); with partial refresh the parent history has memory
(\MT{rem:rg-memory-partial-refresh}), whose first-order cost is $D^{(1)}_N$.

\subsection{The retained-state table and what ``exact'' means}

\MT{tab:rg-retained-state} lists each retained object with its exact operation, the statement that
derives it, and its admission hypotheses, and the chapter says it is a catalog, not a completeness
theorem. Its first row keeps the complete normalized generative, posterior, and recognition laws, and
its last the path law with its memory. Closure is therefore achieved by keeping everything; the other rows
name parts of those objects that a smaller description must not discard, several with witnesses that
they cannot be dropped in general (hyperedges, receiver laws, memory, the marked kernel rather than its
mean, and in \cref{comp:ch12} the graph datum). ``Exact'' means exact as $\sum_ye^{-H(y)}$ is exact, a
lossless change of description, not exact as one-dimensional Ising decimation is, a finite set of
couplings with a closed recursion; no reduced retained set is shown to close. Blocks and channel are
inputs, and must be: at the uniform law on seven outcomes, no rule that sees only the law and its
Fisher metric and treats all outcomes alike can select three two-dimensional blocks
(\MT{prop:hist-law-fisher-agentization-nogo}).

\section{What is proved, and how solid it is}
\label{comp:ch11-scorecard}

\begin{scorecard}
\MT{thm:rg-effective-action} & Coarse action is minus the log of a conditional partition function;
evidence preserved & \SOLID & Any channel; standard in substance\\
\MT{thm:rg-bounded-action-calculus}, \MT{thm:rg-action-lp-contraction} & Analytic, nonexpansive action
increment; cumulant derivatives; $L^p$ and Dobrushin contraction & \SOLID & Whole-action norms;
$\delta=1$ for deterministic block maps\\
\MT{thm:rg-hoeffding-action-isomorphism}, \MT{cor:rg-hoeffding-gauge-covariance} & Couplings and
actions isomorphic; extraction constant $3^{\lvert V\rvert}-1$, sharp; gauge covariant & \SOLID & ANOVA
standard; sharp constant specific\\
\MT{rem:rg-bounded-tier-scope}, \MT{prop:rg-product-equivalence-not-preserved} & Polynomial couplings on
$\R^K$ excluded; product equivalence can fail & \SOLID & Negative; Gaussian $\GL(K)$ couplings are out
of scope\\
\MT{lem:rg-interaction-step}, \MT{prop:rg-interaction-rn-gauge-covariance} & Exact step fixes the origin,
has the tilted derivative, is gauge covariant & \SOLID & Bounded tier; covariance of a.e.\ classes\\
\MT{prop:rg-interaction-native-growth-bound} & Multipliers telescope; norm at most
$3^{\lvert V_n\rvert}-1$ at every depth & \SOLID & Bounded tier; vacuous as a number for large targets\\
\MT{prop:rg-interaction-finite-network-obstruction} & No identification of coupling spaces across a
reducing step; blocking terminates & \SOLID & Negative; every finite network, not only small ones\\
\MT{prop:rg-interaction-residual}, \MT{rem:rg-trivial-invariant-pairs} & Residual identity; independence,
symmetry, and pullback sectors are invariant & \SOLID & No residual estimate; trivial pairs only\\
\MT{rem:rg-two-agent-example} & Decimating one binary agent & \TOY & Two agents; textbook mechanism\\
\MT{prop:rg-hyperedge-generation} & Pairwise closure fails (star), not universally (parity) & \SOLID &
Negative by counterexample; ``generic'' within the star family\\
\MT{prop:rg-bridge-kernels} & Bridge kernels adjoint; no reconstruction of the fine posterior & \SOLID &
Standard in substance\\
\MT{prop:rg-meta-attention-pushforward}, \MT{prop:rg-attention-per-receiver-collapse} & Push event laws;
evidence-weighted versus prior-weighted receivers & \SOLID & Inside two declared label models\\
\MT{thm:rg-strong-lumpability}, \MT{thm:rg-projection-memory}, \MT{thm:rg-memory-filter} & Lumpability;
discrete Nakajima--Zwanzig; exact filter & \SOLID & Standard in substance\\
\MT{thm:rg-memory-projection}, \MT{thm:rg-memory-duration} & Memory-$k$ cost; excess entropy at most
hidden entropy & \SOLID & Finite spaces; stationary start; budget extensive in $N$\\
\MT{tab:rg-retained-state}, read as closure & Exact closure keeps the listed objects & \VAGUE & Rows
proved or declared; sufficiency trivial; no minimality claim\\
\MT{prop:hist-law-fisher-agentization-nogo} & Law and Fisher metric cannot select blocks & \SOLID &
Negative\\
\MT{def:rg-network-scale}, \MT{def:rg-retained-projection} & Channel, blocks, reference, retained
projection are inputs & \CHOICE & Declared, not selected\\
\MT{open:rg-invariant-retained-subspace}, \MT{open:rg-pointwise-parent-dynamics} & Closed truncation;
closed parent dynamics & \OPEN & See \cref{comp:ch11-open}\\
\end{scorecard}

The only toy positive is the two-agent decimation; the other small instances are counterexamples, hence
solid negative results. The mechanisms they display (fielded eliminated agents generate multi-agent
couplings, hidden coordinates produce memory, evidence weighting differs from prior weighting)
plausibly survive at large $N$ but are proved only where stated. The general results hold at every
finite $N$; their only $N$-dependence is in $3^{\lvert V_n\rvert}-1$ and $H(Y_0)$, the quantities a
large-$N$ theory would have to replace with intensive ones.

\begin{hygiene}
The conditions that can fail in the program's models are these. The bounded tier admits finite,
categorical, and compact group-valued states ($\mathrm{SU}(N)$ included) and excludes the quadratic
couplings of $\R^K$-valued agents, so the coupling results say nothing about the Gaussian $\GL(K,\R)$ or
$\GL(K,\mathbb C)$ realization. Product equivalence at both depths fails for channels that impose hard
constraints among coarse agents. The channel and blocks are declared and must not read the recognition
law. Everything is finite-network, with no limit taken. The memory results need finite hidden spaces
and, for the budget, a stationary start. The rest is hygiene (standard Borel spaces, measurable reverse
kernels, almost-everywhere versions, integrability), with one item touching noncompact groups:
pointwise equivariant conditional versions are settled only for finite and compact stabilizers and
weakly continuous posteriors (\MT{open:cg-equivariant-versions}), so for $\GL(K)$ covariance holds for
almost-everywhere classes.
\end{hygiene}

\section{Open ends}
\label{comp:ch11-open}

\MT{open:rg-invariant-retained-subspace} asks for a closed truncation. It is posed for finite agent
spaces with full-support product references, at least two blocks, and deterministic block-product
channels whose block maps read every member of their block (majority or block-sum maps, not
decimation), for product base laws and for all base laws. It is settled by a proper invariant pair,
invariant for every base law in an open set, whose range is not inside the sum of the pullback sector
and any grouping sector, or by a proof that generically every proper invariant pair lies in such a
sum. The exclusion of decimation sets aside the classical exactly closed recursion of the
one-dimensional Ising chain \citep{goldenfeld1992lectures}.

\MT{open:rg-pointwise-parent-dynamics} is the defining problem. The exact parent exists at every
instant as the image of the fine law, so parents do evolve, but only as images of the fine evolution;
the question is whether that evolution is a dynamics of its own. In a differentiable case the
obstruction is the semiconjugacy defect
\[
\delta_t=\partial_t\mathsf F_t+D\mathsf F_t\mathscr V_t-\overline{\mathscr V}_t\circ\mathsf F_t
\]
between the fine field pushed through the moving coarse map and a candidate parent field; for Markov
evolutions the exact condition is semigroup intertwining. Existing statements answer the unrestricted
question both ways (lumpable kernels lump, a parent carrying the full filter reproduces every path law,
the three-cycle admits no first-order parent), so it is posed for a class: switching diffusions at
$c_*$ of the four-cycle type of \cref{comp:ch13} (finitely many source assignments, a quadratic energy
on each assignment's sheet, overdamped diffusion, a rate-one clock redrawing the assignment), with a
parent that keeps the assignment and a projection of the configuration but hides coordinates that the
receiver rows of a non-negligible assignment read. Candidate parents are time-homogeneous Markov
processes of the retained variables, possibly with a regular finite-dimensional filter
(\MT{def:rg-regular-finite-filter}), started from the exact image. It asks whether every member admits
such a parent whose path law agrees in total variation as the regime parameter $\epsilon\to0$, on fixed
windows up to the first expected switch and on whole histories with diverging numbers of switches; a
member for which every such parent stays a bounded distance away settles it negatively. In physics terms: does adiabatic elimination with finitely many auxiliary variables close
the retained dynamics when hidden coordinates drive the switching? The one positive instance is the
four-cycle's inherited parent on one episode window (\MT{prop:rg-formation-endogenous-cycle}), a toy
case. The limit is $\epsilon\to0$ at any finite size, not $N\to\infty$.

Three gaps are not posed as open problems of the main manuscript: no truncation residual is estimated;
the unbounded tier on $\R^K$ exists only in the Gaussian case of Supplement G; and a thermodynamic limit
of the exact step belongs to Supplement B, where fixed-depth and diagonal limits and universality are
open (\SB{open:rg-two-index-limits}).

\begin{verdictbox}{\NARROWER}
The chapter says it determines what one exact step does to the agent network at $c_*$ and which state
must be kept so that the step closes, with an effective action, residual, finite-step multipliers, and
native growth bound; the introduction calls this exact coarse-graining of the agent network. The proofs
deliver those objects: the effective action for every channel, its analytic nonexpansive calculus, an
exact step on couplings with an exact residual identity, a depth-independent derivative bound, exact
memory identities and costs, and witnesses that hyperedges, receiver laws, marked kernels, and memory
cannot be dropped in general. The delivery is narrower than the wording in four respects. Closure is
achieved by keeping the complete laws; no smaller closed set and no truncation with a bounded error are
proved. The coupling machinery holds only in the bounded tier, which excludes the program's Gaussian
$\GL(K)$ couplings. Everything is finite-network, so there is no flow on fixed couplings (proved
impossible across reducing steps) and no thermodynamic limit. And parent dynamics is open. The scope
paragraph says most of this; ``which state must be kept so that the step closes'' is where a physicist
will read more than is there. Much of the content is standard in substance; specific are the typing at
$c_*$, the sharp constant and telescoping bound, the obstruction, the attention-event bookkeeping, and
gauge covariance throughout.
\end{verdictbox}

\chapter{The Multiscale Graph VFE and Closing the Tower}
\label{comp:ch12}

\section{The question and the answer}
\label{comp:ch12-question}

The technical chapter (\MT{ch:full-graph-meta-agent-vfe}) asks how to place everything exact
coarse-graining produces (adaptive attention event laws and links, random partitions, and the influence
of parents on their children) into one variational free energy across depths, and how a tower of depths
ends. Its answer is one normalized generative law and one recognition law over the whole tower, with
free energy
\[
\Fenergy^{\mathrm{ext}}_{\mathrm{tower}}[\mathbb Q_\lambda;X_*,o]
=-\log p_\theta(o\given X_*)
+\KL\bigl(\mathbb Q_\lambda(\cdot\given o,X_*)\Vert\boldsymbol\Pi_{\theta,o,X_*}\bigr)
\]
(\MT{def:full-graph-tower-vfe}); everything else in the chapter is exact bookkeeping of this object or
is labeled as something else. For a tower built bottom-up by coarse-graining channels, both laws factor
top-down automatically, each depth conditioned on the complete next-depth state (coordinates and graph
datum) and the partition mark (\MT{lem:full-graph-adaptive-reversal}); dropping the graph datum costs a
sum of mutual informations (\MT{prop:full-graph-dropped-datum}); and the free energy splits into
observation, top-prior, partition, graph, and parent--child relative entropies
(\MT{thm:full-graph-vfe-decomposition}). Partitions may be random, from a finite catalog under two
inequivalent normalizations (\MT{prop:rg-geometry-weights}). The tower cannot close on itself: a finite
acyclic model has a parentless latent with a declared prior (\MT{prop:obs-declared-root-unavoidable}).
What remains of the idea that the agents constitute their own top lives at three loci (the inference
loop, learning of the top prior, and delayed feedback across episodes), each an evaluation of fixed
normalized joints, and delayed feedback with a declared update rule is itself the free energy of one
lifted joint (\MT{prop:full-graph-episode-lifted-joint}). Whether descent of this free energy produces a
persistent hierarchy is open (\MT{open:full-graph-hierarchy}).

Do parents act back on their children in a proved sense? In four senses of different strength. In a
declared top-down hierarchy the downward kernel reads the parent coordinates, so the parent enters every
child's prior by construction, a modeling choice; parent and child optimization are then coupled
exactly when a mixed derivative is nonzero, which the chapter proves only as a criterion in a binary
example (a kernel odds ratio different from one), declining any family-wide claim without a proof in
that family (\MT{prop:full-graph-downward-coupling}). In an adaptive tower the downward kernel is the
Bayes reversal of the bottom-up coarse-graining, and the chapter declines to read it causally. Through
the inference loop, each child's effective prior at a mean-field fixed point is formed from the parent's
factor, itself formed from all children (\MT{rem:full-graph-inference-loop}), a standard fact. Across
episodes, the parent inferred at one episode sets the children's prior at the next inside one declared
joint (\MT{prop:full-graph-episode-joint}). Not proved is a parent with its own dynamics acting back over
time: that is \MT{open:rg-pointwise-parent-dynamics} and the third result required by
\MT{open:full-graph-hierarchy}.

\section{The picture}
\label{comp:ch12-picture}

Picture depths $\ell=0,\ldots,L$: agents at depth $0$, meta-agents at depth $1$, a declared top at depth
$L$. The state at each depth includes, besides coordinates, the graph datum: who attends to whom in both
channels, the links on every edge copy, and the root and holonomy marks of each block. Going up, a
structural channel takes the complete depth-$\ell$ state and produces a partition mark $R_\ell$ and the
depth-$(\ell+1)$ state. Read downward, the same joint law is a Markov chain in decreasing depth whose
transitions are the conditional laws of finer given coarser states. The free energy of the whole stack
is one Gibbs variational functional, and its per-depth decomposition is the chain rule.

\begin{analogy}{exact}
Bayes reversal of an adaptive tower is time reversal of a Markov chain with depth as time: the chain
$Y_0\to(R_0,Y_1)\to(R_1,Y_2)\to\cdots$ read backward is Markov, and its downward transition is the
reverse block-spin kernel of \cref{comp:ch11}. No downward causation follows, as time reversal of a
diffusion does not make the future cause the past.
\end{analogy}

\begin{analogy}{exact}
The two normalizations of a partition catalog are the two ways to combine Gibbs measures on several
sectors: a mixture of separately normalized measures with weights fixed by hand (hierarchical), or one
Gibbs measure on the disjoint union, in which each sector weighs its a priori weight times its partition
function, $\mathsf a_\gamma e^{-F_\gamma}$ (global). The first fixes the probabilities of the phases; the
second lets their free energies take part.
\end{analogy}

\begin{analogy}{exact}
The coarse links of \MT{def:full-graph-coarse-links} are maximal-tree gauge fixing inside each block
\citep{Creutz1977}: tree links go to the identity, the remaining internal links carry the block's Wilson
loops as self-loops of the coarse vertex, and boundary links connect block roots. Unlike smearing of
gauge links, nothing is averaged or projected, and every parallel copy is kept.
\end{analogy}

\section{Objects and notation}
\label{comp:ch12-notation}

\begin{decoder}
$Z_\ell$ & Full coordinates at depth $\ell$ (beliefs and models) & Block variables\\
$\mathfrak N_\ell=(\eta_\ell,\Theta^{b,\ell},\Theta^{m,\ell},\mathfrak h_\ell)$ & Graph datum: event laws,
links, root and holonomy marks & Coarse lattice with gauge field and bond occupations\\
$Y_\ell=(Z_\ell,\mathfrak N_\ell)$, $\mathbf Y$ & Complete state at depth $\ell$; all depths together &
Block configuration; multiscale configuration\\
$R_\ell$, $\gamma$ & Structural mark (realized partition or membership pattern) & Random blocking
pattern\\
$C^{\mathrm{struct}}_{\ell+1\leftarrow\ell}$ & One kernel from $Y_\ell$ to $(R_\ell,Y_{\ell+1})$ applied to
all three laws & Configuration-dependent block transformation\\
adaptive tower & Tower generated bottom-up by structural channels & Iterated blocking\\
$K^\ell_\downarrow$ & Law of depth-$\ell$ coordinates given the next depth & Prolongation; reverse
block-spin kernel\\
$\Fenergy^{\mathrm{att},x}_{i,\ell}$ & Row free energy: mean edge divergence plus $\tau$ times relative
entropy to the prior & Helmholtz free energy of a receiver's source choice\\
$\mathcal J_{\mathrm{comp}}$ & Energies appended without a normalizer & Lyapunov function, not a free
energy\\
$\mathfrak G$, $\mathsf a_\gamma$, $\mathcal Z_\gamma$ & Partition catalog; label weights; branch
partition functions & Sectors, their weights and partition functions\\
$\mathfrak m_\ell$, $\mathfrak M_\ell$ & Soft membership of a vertex; joint membership of an edge's
endpoints & Fuzzy block assignment\\
declared root, $\mathbb P_L$ & Top prior supplied by no other latent & Boundary condition at the top\\
episode $n$, $\theta_n$, $\mathcal U_n$ & Learning episode, parameter, declared update & Sweep of a
learning loop\\
latent route; parameter route & One joint across episodes; separate joints linked by updates & Hidden
Markov model over episodes; learning by updates\\
lifted joint, $\widetilde o_n$ & One joint with parameters and record copies latent; latent record copy
& State augmentation in filtering\\
$\operatorname{TC}$, $\epsilon_{\mathrm e}$, $\mathrm T$ & Total correlation across episodes; aggregate
update step; descent temperature & Multi-information; learning rate; temperature\\
\end{decoder}

Here $Y_\ell$ includes the graph datum and $Z_\ell$ is the coordinate block at depth $\ell$, unlike the
coarse sample $z$ of \cref{comp:ch11}. A learning episode (index $n$) is not a coherence episode of
formation (\cref{comp:ch13}). The fraktur $\mathfrak m_\ell$ is a membership kernel, unrelated to the
evidence measure $\mathsf m_\ell$ of \cref{comp:ch11}, and $R_\ell$ is a partition mark, not a rotation.
``Top-down'' names the direction of a generative arrow and, in an adaptive tower, carries no causal
meaning. ``Closing the tower'' means ending it with a declared root. ``Participation'' is interpretive
vocabulary, not a mathematical property.

\section{The formulas, explained}
\label{comp:ch12-formulas}

\subsection{One free energy and the structural channel}

The tower free energy is the evidence identity of \cref{comp:ch05} on a latent space that contains every
depth; with $L=0$ it is the population free energy of one network. The commitment, and the reason it is
a choice rather than a theorem, is that one normalized generative law covers all depths: independently
chosen per-depth energies added together do not define a free energy, and the chapter calls such a sum a
composite potential.

At depth $\ell$ the closed scale datum holds the full coordinates and the graph datum on a finite
multigraph, the support of the source prior at depth $0$ and the multigraph induced by coarse links
deeper down (\MT{def:full-graph-scale-datum}). A structural step is one recognition-independent channel
$C^{\mathrm{struct}}_{\ell+1\leftarrow\ell}:\mathsf Y_\ell\rightsquigarrow\mathsf R_\ell\times\mathsf
Y_{\ell+1}$, applied alike to the generative marginal, the posterior, and the recognition law, whose
laws of $R_\ell$ may therefore differ (\MT{def:full-graph-structural-channel}). Kadanoff blocks are a
constant mark with a deterministic block-product map. Random blocks are a mark drawn by the channel,
possibly from the fine configuration, which makes the step adaptive. What the blocking may not read is a
law: choosing blocks by inspecting the posterior or recognition law is excluded, since that is not one
channel applied to three laws. Soft blocks are membership kernels $\mathfrak m_\ell(I\given i,R_\ell)$, a
hard partition being the deterministic case, with a joint endpoint kernel $\mathfrak M_\ell$ for
correlated assignment of an edge's two ends (\MT{def:full-graph-membership-kernel}). Under independent
draws the coarse event law is
\[
\eta^{x,\ell+1}_{IJ}(y,\gamma)
=\sum_{i\ne j}\eta^{x,\ell}_{ij}(y)\mathfrak m_\ell(I\given i,\gamma)\mathfrak m_\ell(J\given j,\gamma)
+\delta_{IJ}\sum_i\eta^{x,\ell}_{ii}(y)\mathfrak m_\ell(I\given i,\gamma),
\]
the loop term using one draw for both ends of a self-edge (\MT{prop:full-graph-event-pushforward}). The
random $\eta_{\ell+1}$ is part of the state, not its mean, and unit-weight membership in several parents
replicates mass rather than splitting it. Nested assignments compose by Chapman--Kolmogorov, and a
multi-depth channel is exact only if the state carries every mark the next kernel reads
(\MT{def:full-graph-composition}).

\subsection{Adaptive reversal and the cost of dropping the graph datum}

The standing hypothesis \MT{hyp:full-graph-depth-markov} factors the generative law from the top,
\[
\begin{aligned}
\mathbb P_\theta(do,d\mathbf Y\given X_*)
={}&\mathbb P^{\mathrm{obs}}_\theta(do\given Z_0,X_*)\mathbb P_L(dY_L\given X_*)\\
&\times\prod_{\ell=0}^{L-1}\mathbb P^\ell_R(dR_\ell\given Y_{\ell+1})
\mathbb P^\ell_{\mathfrak N}(d\mathfrak N_\ell\given Y_{\ell+1},R_\ell)
K^\ell_\downarrow(dZ_\ell\given Y_{\ell+1},R_\ell,\mathfrak N_\ell),
\end{aligned}
\]
with $X_*$ suppressed in the depth factors, and recognition likewise. Each depth reads the complete next
depth and nothing deeper: a Markov property in scale, not mean field, since dependence within and between
adjacent depths is arbitrary. For a tower generated bottom-up by structural channels it is automatic
(\MT{lem:full-graph-adaptive-reversal}): deeper variables are produced from $Y_{\ell+1}$ by kernels that
see nothing finer, and successive disintegration gives the top-down factors, which are reversal (bridge)
kernels rather than independent choices. When the parent is a function of the child and the mark, the
downward kernel is supported on compatible child configurations, the constrained sum of a Kadanoff block
read backward.

If the depth-$\ell$ factors read only $(Z_{\ell+1},R_\ell)$, the reduced product differs from the tower
law by
\[
\KL\bigl(\mathbb P^{\mathrm{tow}}\Vert\mathbb P^{\mathrm{red}}\bigr)
=\sum_{\ell=0}^{L-1}I_{\mathbb P^{\mathrm{tow}}}\bigl(Y_\ell;\mathfrak N_{\ell+1}\given Z_{\ell+1},R_\ell\bigr)
\]
(\MT{prop:full-graph-dropped-datum}): what the coarse graph datum says about the fine state beyond the
coarse coordinates and the partition. The chapter's four-vertex witness, where fine event laws with equal
coarse coordinates push to a boundary event or a meta-self event, costs $\log2$. The identity is the
chain rule; its content is that coarse attention structure is part of the coarse state.

\subsection{The hierarchical decomposition and the row free energies}

Under the depth-Markov factorization, with common domination and integrability,
\[
\begin{aligned}
\Fenergy_{\mathrm{tower}}
={}&-\E_{\mathbb Q_\lambda}\log p^{\mathrm{obs}}_\theta(o\given Z_0,X_*)
+\KL(\mathbb Q_L\Vert\mathbb P_L)\\
&+\sum_{\ell=0}^{L-1}\E_{\mathbb Q_\lambda}\bigl[
\KL(\mathbb Q^\ell_R\Vert\mathbb P^\ell_R)
+\KL(\mathbb Q^\ell_{\mathfrak N}\Vert\mathbb P^\ell_{\mathfrak N})
+\KL(\mathbb Q^\ell_Z\Vert K^\ell_\downarrow)\bigr]
\end{aligned}
\]
(\MT{thm:full-graph-vfe-decomposition}): the data misfit at the finest depth, the top's recognition law
against its declared prior, and at each depth the costs of the partition mark, the graph datum, and the
coordinates against their parent-conditioned priors, averaged under recognition. With one depth it is
misfit plus prior relative entropy. It is the measure-level chain rule, standard in substance; its value
is that every candidate hierarchy has exact accounting.

With transported divergences $D^{x,\ell}_e$, copy priors $\pi^{x,\ell}_i$, and temperatures
$\tau^x_{i,\ell}$, the row free energy
$\Fenergy^{\mathrm{att},x}_{i,\ell}=\sum_e\beta^{x,\ell}_eD^{x,\ell}_e+\tau^x_{i,\ell}\KL(\beta^{x,\ell}_i\Vert\pi^{x,\ell}_i)$
(\MT{def:full-graph-transported-divergences}) is uniquely minimized by the Gibbs row
$\beta^*_e\propto\pi_ee^{-D_e/\tau}$ with minimum $-\tau\log\mathcal Z$, infinite-energy edges included
(\MT{prop:full-graph-gibbs-rows}): the Gibbs variational principle for $U-\tau S$ with the prior as
degeneracy. With $\tau\ne1$ this energy-unit form is not an independently weighted sector of one
evidence bound. Appending such row energies, partition terms, cross-scale terms, and a link regularizer
to a fixed population free energy gives the composite potential of
\MT{def:full-graph-composite-potential}, a descent potential with no evidence interpretation; the same
terms from a normalized Gibbs prior inside the joint belong to the tower free energy instead.

\subsection{Random partitions and the two normalizations}

A finite catalog $\mathfrak G$ gives each label $\gamma$ a carrier, a reference $\nu_\gamma$, a raw factor
$\widetilde p_\gamma$, a weight $\mathsf a_\gamma$, and a likelihood, with branch partition function
$\mathcal Z_\gamma=\sum_y\nu_\gamma(y)\widetilde p_\gamma(y)$ (\MT{def:rg-geometry-ensemble}):
\[
\mathbb P_{\mathrm h}(\gamma,y,o)=\bar{\mathsf a}_\gamma p_\gamma(y)L_\gamma(o\given y),
\qquad
\mathbb P_{\mathrm f}(\gamma,y,o)=\frac{\mathsf a_\gamma\nu_\gamma(y)\widetilde p_\gamma(y)L_\gamma(o\given y)}{\sum_\delta\mathsf a_\delta\mathcal Z_\delta}.
\]
Both are normalized with the same within-branch laws, but the label weight is $\bar{\mathsf a}_\gamma$
hierarchically and proportional to $\mathsf a_\gamma\mathcal Z_\gamma$ globally; the joints coincide
exactly when all $\mathcal Z_\gamma$ agree, rescaling one branch's raw factor changes only the global
joint, and the two free energies differ only in normalizers (\MT{prop:rg-geometry-weights}). A catalog
creates no edge. The partition term of \MT{def:full-graph-partition-factor} enters the free energy only
when its Gibbs factor is part of the generative law, and no canonical partition selector follows from the
free energy (\MT{open:cg-natural-selector}).

\subsection{Coarse links}

Under a hard partition with one root and spanning tree per block, every fine edge copy survives between
the blocks of its endpoints, with the dressed link
$\Theta^{x,\ell+1}_e=\Theta^{x,\ell}_{I\leftarrow i}\Theta^{x,\ell}_e(\Theta^{x,\ell}_{J\leftarrow j})^{-1}$
(\MT{def:full-graph-coarse-links}). By \MT{prop:full-graph-coarse-links}, no link content is lost: internal
tree copies have identity links, internal non-tree copies carry the based holonomy of their fundamental
loops, closed-walk products reproduce the block holonomy, and coarse divergences and rows are gauge
invariant. A loop divergence vanishes on a non-tree copy exactly when the parent law is fixed by that
holonomy. Once holonomy is nontrivial the trees are part of the coarse channel, not a gauge choice,
although the divergences do not depend on them for holonomy-stabilized parent laws. Overlapping
membership is excluded, since an edge would need a root transport into every block containing an
endpoint.

\subsection{How a parent enters its children's prior}

In a declared top-down hierarchy whose downward kernel factors over parents and reads the parent
coordinates, the cross-scale term of the free energy is
\[
\Fenergy^\ell_{\mathrm{cross}}
=\E_{\mathbb Q_\lambda}\KL\bigl(\mathbb Q^\ell_Z(\cdot\given Y_{\ell+1},R_\ell,\mathfrak N_\ell,o,X_*)\Vert
K^\ell_\downarrow(\cdot\given Y_{\ell+1},R_\ell,\mathfrak N_\ell,X_*)\bigr)
\]
(\MT{prop:full-graph-downward-coupling}), and the parent can also couple to child marks through the
partition and graph factors. A kernel that ignores the parent contributes no parent--child
cross-derivative through this factor, and normalization alone implies no coupling. When the kernel reads
the parent and the recognition parameterization transmits the dependence, the mixed derivative can be
nonzero, so optimization is coupled both ways although the generative arrow points down. For a binary
parent and child with product recognition $q_b(1)=u$, $q_y(1)=v$, the mixed derivative is
$-\log[K_\downarrow(1\given1)K_\downarrow(0\given0)/(K_\downarrow(0\given1)K_\downarrow(1\given0))]$,
zero exactly at unit odds ratio. That criterion in a two-variable example is the whole of the proved
bidirectional coupling.

\subsection{Closing the tower and the three loci}

In a finite directed acyclic model where every record has a latent parent, a topological order starts at
a parentless latent whose factor depends only on parameters and structural data
(\MT{prop:obs-declared-root-unavoidable}). The proof is one line of graph theory and the result an
established negative: within one finite acyclic joint the top cannot be constituted by the agents it
organizes. Same-time reciprocal structure among latents is admissible only as a globally normalized
undirected model with its partition function (\MT{sec:gen-undirected}, \MT{prop:gen-gibbs-counterexample});
an energy appended without its normalizer is not an evidence bound.

The first locus is the inference loop (\MT{rem:full-graph-inference-loop}): mean-field coordinate updates
of parent and child read each other \citep[Sec.~10.1.1]{Bishop2006}, and at a fixed point a child's
effective prior is formed from the population's beliefs through a latent inside the joint. This is
mean-field self-consistency, as when each spin of a mean-field magnet sees a field set by its neighbors,
whose fields include its own; it is a schedule against one fixed joint, not participation in time. The
second is learning the top prior (\MT{prop:full-graph-aggregate-prior}): with independent episodes
sharing a top-prior parameter, the optimum over an unrestricted family is the aggregate recognition law,
the variational form of empirical Bayes \citep{Robbins1956} and the aggregate-posterior prior of
\citet{HoffmanJohnson2016} and \citet{TomczakWelling2018}.

The third is delayed feedback. In the latent route (\MT{def:full-graph-episode-joint}) one joint spans
episodes, and the parent read out of episode $n$ parameterizes episode $n+1$, so
\[
\mathbb P^{\mathrm{ep}}\bigl(\mathbf Y^{(n+1)}\in\cdot\given o_{0:n}\bigr)
=\int\mathbb P^{(n+1\leftarrow n)}(\cdot\given z,X)
\mathbb P^{\mathrm{ep}}\bigl(Z^{(n)}_{\mathrm{par}}\in dz\given o_{0:n}\bigr)
\]
(\MT{prop:full-graph-episode-joint}): the children's next prior is the filtered parent pushed through the
transition. The joint is acyclic, so it still needs a declared root at episode $0$. In the parameter
route each episode has its own joint and the parameter is updated by a declared map of the selected
recognition coupling. \MT{prop:full-graph-episode-lifted-joint} shows that such a route is one joint:
make the parameters and latent record copies $\widetilde o_n$ latent states with Dirac transitions and
evaluate the update at the copy, not at the recognition law being scored. The lifted joint's evidence is
the product of the episode evidences, its posterior the product of episode posteriors, and
\[
\Fenergy^{\mathrm{lift}}[\mathbb Q]
=\sum_{n=0}^{n_1}\Bigl[-\log p_{\theta_n,V}(o_n\given X)
+\KL\bigl(\mathbb Q^{y}_n\Vert\boldsymbol\Pi_{\theta_n,V,o_n,X}\bigr)\Bigr]
+\operatorname{TC}(\mathbb Q^{y}),
\]
the sum of the episode bounds plus a cross-episode total correlation that vanishes at product laws. In
physics terms this is state augmentation: a learning loop becomes filtering once the parameters are
latent. The point is typing: no generative factor reads the recognition law being scored, so the typing
prohibition of \cref{comp:ch04} (\MT{req:gen-typing-prohibition}) holds.

\section{What is proved, and how solid it is}
\label{comp:ch12-scorecard}

\begin{scorecard}
\MT{def:full-graph-tower-vfe} & One generative and one recognition law over the whole tower & \CHOICE &
The modeling commitment; the rest is accounting\\
\MT{def:full-graph-structural-channel}, \MT{def:full-graph-membership-kernel} & One structural channel
for three laws; blocking reads states, not laws; soft membership & \CHOICE & Covers Kadanoff, random,
adaptive, and soft blocks\\
\MT{lem:full-graph-adaptive-reversal} & Adaptive towers factor top-down; downward factors are reversal
kernels & \SOLID & Standard in substance\\
\MT{prop:full-graph-dropped-datum} & Dropping the graph datum costs a sum of mutual informations &
\SOLID & Chain rule; $\log2$ witness illustrative\\
\MT{thm:full-graph-vfe-decomposition} & Five-term decomposition of the tower free energy & \SOLID &
Standard in substance; depth-Markov form\\
\MT{prop:full-graph-gibbs-rows} & Gibbs rows minimize row free energies & \SOLID & Gibbs variational
principle; $\tau\ne1$ not a sector of one bound\\
\MT{def:full-graph-composite-potential} & Appended energies give a descent potential & \CHOICE & No
evidence interpretation, as stated\\
\MT{prop:rg-geometry-weights} & The two normalizations differ by branch partition functions & \SOLID &
Finite catalog and carriers\\
\MT{prop:full-graph-event-pushforward}, \MT{def:full-graph-composition} & Event pushforward under soft
membership; composition & \SOLID & Product form needs independent draws\\
\MT{prop:full-graph-coarse-links} & Coarse links keep all holonomy; trees belong to the channel & \SOLID
& Hard partitions only\\
\MT{prop:full-graph-downward-coupling}, general part & No parent dependence, no cross-derivative;
normalization gives no coupling & \SOLID & Declared top-down hierarchy\\
\MT{prop:full-graph-downward-coupling}, criterion & Bidirectional coupling iff odds ratio $\ne1$ & \TOY &
Binary parent and child; no family theorem\\
\MT{prop:obs-declared-root-unavoidable} & A finite acyclic tower has a declared root & \SOLID & Negative;
every record needs a latent parent\\
\MT{rem:full-graph-inference-loop} & Parent and child updates read each other & \SOLID & Standard\\
\MT{prop:full-graph-aggregate-prior} & Optimal top prior is the aggregate recognition law & \SOLID &
Empirical Bayes in substance\\
\MT{prop:full-graph-episode-joint}, \MT{prop:full-graph-episode-lifted-joint} & Delayed feedback is the
free energy of one joint & \SOLID & Declared measurable selection and update\\
\MT{sec:obs-participatory}, read as participation & The loci realize what a cycle was reaching for &
\VAGUE & Interpretation; each identity proved\\
\MT{open:full-graph-hierarchy}, \MT{open:full-graph-episode-accounting} & Persistent hierarchy; feedback
by labels alone & \OPEN & See \cref{comp:ch12-open}\\
\end{scorecard}

Toy content is confined to the binary odds-ratio criterion and the four-vertex $\log2$ witness; only the
first stands in for a general claim, which the chapter explicitly does not make.

\begin{hygiene}
Four conditions matter. The depth-Markov form is automatic for adaptive towers but must be built into a
declared hierarchy. The structural channel must not read a law, which excludes posterior-dependent
blocking. Coarse links are constructed for hard partitions only. And no generative factor may read the
recognition law being scored, which is why the lifted joint evaluates its update at a latent record
copy. A fifth concerns reading, not truth: energy-unit row free energies with $\tau\ne1$ and the
composite potential are not evidence bounds. The rest is hygiene: standard Borel spaces and regular
conditional versions, common domination and integrability, extended-real conventions, joint measurability
of selector and update, and discrete records for the displayed episode decomposition.
\end{hygiene}

\section{Open ends}
\label{comp:ch12-open}

\MT{open:full-graph-hierarchy} asks for the hierarchy the program is after: a normalized partition law
and parent--child kernels at $c_*$ such that descent of the tower free energy produces persistent nested
blocks, with pushforwards and link lifts that compose, a controlled closure residual, and non-flat
holonomy kept rather than forced flat. What holds is the exact accounting, the fast--slow typing of
\cref{comp:ch08}, nested formation under a supplied hierarchy of interaction strengths
(\MT{thm:finite-depth-adaptive-formation}), and one endogenous level on the four-cycle
(\MT{prop:rg-formation-endogenous-cycle}). Unproved is that descent produces a second level that
depends on the first. The problem is posed for models with a partition catalog inside both laws,
recognition coordinates descending by reversible diffusion at temperature $\mathrm T$ with rate-one mark
redraws, and $\mathrm T$-independent data, as $\mathrm T\to0$. A positive answer needs three results: a
two-level lineage formed with probability bounded below whose second level is recursive in the
manuscript's strict sense (it depends on the first by a margin over counterfactual handoffs and is not a
supplied-scale sequence) and is a coherence episode (held by a restoring interaction against thermal
noise, on horizons short compared with its stationary recurrence or with probability exceeding its
stationary value by a margin); an intrinsic partition selector; and a separated slow manifold on which
coarser variables feed back consistently into finer kernels. A negative answer is a proof that no second
level produced by descent in this class meets these conditions. The recursion and coherence rules are
decoded in \cref{comp:ch13}. The limit is zero temperature on networks of any finite size, not
$N\to\infty$.

\MT{open:full-graph-episode-accounting} asks whether delayed feedback can be carried by the labels alone.
In a label-sufficient latent route the episode model is that of a parameter route with the aggregate
update
$\theta_{n+1}=(1-\epsilon_{\mathrm e})\theta_n+\epsilon_{\mathrm e}\boldsymbol\Pi^{\mathrm{lab}}_{\theta_n,V,o_n,X}$,
and the parent read-out sees only the label and a preference coordinate. Does some such route's
filtered label law follow the aggregate update to $o(\epsilon_{\mathrm e})$ per episode, uniformly over
histories and horizons of order $1/\epsilon_{\mathrm e}$, or does every route miss by order
$\epsilon_{\mathrm e}$? The records are restricted either to include an uninformative value while moving
the label posterior by an amount bounded below on average, or to have generic likelihood vectors (no
nonzero polynomial vanishes on their cone). In physics terms: can a stochastic-approximation learning
rule be represented, to first order in its step, as filtering in a fixed latent model whose memory is a
label and one preference coordinate? No natural partition selector is known either
(\MT{open:cg-natural-selector}), and parent dynamics (\cref{comp:ch11-open}) is what would make the
static and episodic senses of parents acting back dynamical.

\begin{verdictbox}{\DELIVERED}
The chapter says it places adaptive event laws and links, random partitions, and the influence of parents
on their children inside one variational free energy, factors that free energy for adaptive towers, and
closes the tower with a declared root, with participation at three loci. It does exactly that: one Gibbs
variational functional for the tower, automatic top-down factorization, an identity for the five-term
decomposition, an exact cost for dropping the graph datum, the two partition normalizations
distinguished exactly, coarse links that keep all holonomy, a proof that a declared root is unavoidable,
and a proof that delayed feedback with a declared update is the free energy of one joint. Two remarks keep
the verdict honest. What is delivered is accounting and typing, largely standard in substance (chain
rule, Markov-chain reversal, Gibbs variational principle, mean-field coordinate ascent, empirical Bayes,
state augmentation); nothing is proved about what descent of this free energy does, and the hierarchy it
would have to produce is \MT{open:full-graph-hierarchy}. And ``the influence of parents on their
children'' is delivered as a typed possibility: in a declared hierarchy the parent enters the children's
prior by construction, bidirectional coupling is proved only as a binary criterion, and in adaptive
towers the downward kernel is a Bayes reversal the chapter declines to read causally. ``Closing the
tower'' means closing it by declaration; the chapter proves that a finite acyclic tower cannot close
itself.
\end{verdictbox}
